% =====================================================================
% Lab Séance 2 — Analyse algorithmique
% Structures de Données et Algorithmes Avancés — M1 — ESP/UCAD
% Dr. El Hadji Bassirou TOURÉ — 2025–2026
% Compilation : pdflatex (2 passes)
% =====================================================================
\documentclass[11pt,a4paper]{article}

\usepackage[utf8]{inputenc}
\usepackage[T1]{fontenc}
\usepackage{lmodern}
\usepackage[french]{babel}

\usepackage[margin=2.5cm]{geometry}
\usepackage{amsmath,amssymb,amsthm}
\usepackage{graphicx}
\graphicspath{{figures/}}
\usepackage{xcolor}
\usepackage{booktabs}
\usepackage{array}
\usepackage{enumitem}
\usepackage{fancyhdr}
\usepackage{longtable}
\usepackage{multirow}
\usepackage{pifont}
\usepackage{listings}
\usepackage{tcolorbox}
\tcbuselibrary{theorems,breakable,skins}
\usepackage{hyperref}

\hypersetup{colorlinks=true,linkcolor=blue!60!black,urlcolor=blue!60!black,
            citecolor=blue!60!black}

\newcommand{\cmark}{\ding{51}}
\newcommand{\xmark}{\ding{55}}
\newcommand{\bigO}{\mathcal{O}}
\newcommand{\bigOmega}{\Omega}
\newcommand{\bigTheta}{\Theta}
\newcommand{\lig}[1]{\rule{#1}{0.4pt}}

% ===== Boîtes =====
\newtcolorbox{objectifbox}[1][]{colback=blue!5,colframe=blue!50,
    fonttitle=\bfseries,title=#1,breakable}
\newtcolorbox{infobox}[1][]{colback=green!5,colframe=green!50!black,
    fonttitle=\bfseries,title=#1,breakable}
\newtcolorbox{attentionbox}[1][]{colback=orange!5,colframe=orange!70,
    fonttitle=\bfseries,title=#1,breakable}
\newtcolorbox{retenirbox}[1][]{colback=yellow!10,colframe=yellow!50!black,
    fonttitle=\bfseries,title=#1,breakable}
\newtcolorbox{prereqbox}[1][]{colback=cyan!5,colframe=cyan!50!black,
    fonttitle=\bfseries,title=#1,breakable}
\newtcolorbox{theorembox}[1][]{colback=purple!5,colframe=purple!50,
    fonttitle=\bfseries,title=#1,breakable}
\newtcolorbox{appliboxQ}[1][]{colback=gray!8,colframe=gray!60,
    fonttitle=\bfseries,title=#1,breakable}
\newtcolorbox{appliboxR}[1][]{colback=green!3,colframe=green!40!black,
    fonttitle=\bfseries\itshape,title=#1,breakable}
\newtcolorbox{exercicebox}[1][]{colback=white,colframe=gray!70,
    fonttitle=\bfseries,title=#1,breakable}
\newtcolorbox{lienbox}[1][]{colback=blue!3,colframe=blue!70!black,
    fonttitle=\bfseries\itshape,title=#1,breakable,
    before upper={\itshape}}

% ===== Code =====
\lstset{
  language=Java,
  basicstyle=\ttfamily\small,
  backgroundcolor=\color{gray!8},
  frame=single,
  framerule=0.5pt,
  rulecolor=\color{gray!40},
  keywordstyle=\color{blue!70!black}\bfseries,
  commentstyle=\color{green!50!black}\itshape,
  stringstyle=\color{red!60!black},
  numbers=left,
  numberstyle=\tiny\color{gray},
  breaklines=true,
  breakatwhitespace=true,
  showstringspaces=false,
  tabsize=4,
  extendedchars=true,
  xleftmargin=4pt,
  xrightmargin=4pt,
  aboveskip=8pt,
  belowskip=8pt,
  literate={é}{{\'e}}1 {è}{{\`e}}1 {ê}{{\^e}}1 {à}{{\`a}}1 {ù}{{\`u}}1
           {ô}{{\^o}}1 {î}{{\^i}}1 {ç}{{\c{c}}}1 {û}{{\^u}}1 {â}{{\^a}}1
           {ë}{{\"e}}1 {ï}{{\"i}}1 {É}{{\'E}}1 {È}{{\`E}}1 {À}{{\`A}}1
           {Ç}{{\c{C}}}1 {Ê}{{\^E}}1 {œ}{{\oe}}1 {°}{{$^\circ$}}1
           {«}{{\guillemotleft}}1 {»}{{\guillemotright}}1
}
\lstdefinestyle{bash}{language=bash,keywordstyle=\color{black},
  commentstyle=\color{gray!70!black}\itshape,numbers=none}

% ===== En-têtes =====
\pagestyle{fancy}
\fancyhf{}
\fancyhead[L]{\small Lab S02 : Analyse algorithmique}
\fancyhead[R]{\small ESP/UCAD --- M1}
\fancyfoot[L]{\small Dr. El Hadji Bassirou TOURÉ}
\fancyfoot[C]{\thepage}
\fancyfoot[R]{\small 2025--2026}
\renewcommand{\headrulewidth}{0.4pt}
\renewcommand{\footrulewidth}{0.4pt}

\title{%
\vspace{-1cm}
\textbf{\LARGE Structures de Données et}\\
\textbf{\LARGE Algorithmes Avancés}\\[0.8cm]
\textbf{\huge Lab S02 --- Analyse algorithmique}\\[0.3cm]
\textbf{\Large Compter, mesurer, confronter}\\[0.5cm]
\large Master 1 --- 2025--2026}
\author{Dr. El Hadji Bassirou TOURÉ\\
Département Génie Informatique\\
Université Cheikh Anta Diop de Dakar}
\date{}

\begin{document}
\maketitle
\thispagestyle{empty}

\begin{prereqbox}[Avant de commencer]
\textbf{Prérequis :} TP-1 complété, CM~S02 lu, Lab~S01 terminé.\\[4pt]
\textbf{Organisation :} deux parties. \emph{Partie I} --- vous instrumentez
du code pour compter les opérations réellement exécutées, puis vous mesurez
des temps d'exécution. \emph{Partie II} --- vous établissez par le calcul
les formules que vos mesures auront confirmées.\\[4pt]
\textbf{Ce Lab n'est pas noté.}\\[4pt]
\textbf{Livrable :} le dépôt Git du projet \texttt{lab-s02} et le rapport
individuel décrit en fin de document.
\end{prereqbox}

\begin{objectifbox}[Ce que vous saurez faire à la fin de ce Lab]
\begin{enumerate}[nosep]
    \item Instrumenter un fragment de code pour compter ses opérations
          élémentaires selon la convention du cours.
    \item Vérifier expérimentalement une formule de comptage, et localiser
          l'écart quand elle est fausse.
    \item Mesurer un temps d'exécution avec \texttt{System.nanoTime()} en
          évitant les pièges de mesure les plus courants.
    \item Confronter une courbe mesurée à une prédiction asymptotique, et
          interpréter les écarts.
    \item Prouver une borne $\bigO$ en exhibant $c$ et $n_0$.
    \item Résoudre une récurrence de diviser-pour-régner.
\end{enumerate}
\end{objectifbox}

\tableofcontents
\newpage

% #####################################################################
\part{Pratique : instrumentation et mesure}
% #####################################################################

\section{Environnement de travail}

\begin{lstlisting}[style=bash,caption={Rappel : lancer le conteneur et créer le projet}]
# Windows (Git Bash)
MSYS_NO_PATHCONV=1 docker run -it -v "$(pwd):/workspace" elbachirtoure/dsa-java-m1:v1.1
# macOS / Linux
docker run -it -v "$(pwd):/workspace" elbachirtoure/dsa-java-m1:v1.1

cd /workspace
new-project lab-s02
cd lab-s02
gradle build --no-daemon
\end{lstlisting}

\begin{attentionbox}[Deux règles héritées du TP-1]
La classe principale s'appelle \texttt{Main}, et toutes les commandes
Gradle prennent \texttt{--no-daemon}.
\end{attentionbox}

% =====================================================================
\section{Étape 1 --- Compter pour de vrai}

\begin{lienbox}[Ce qu'on matérialise du CM]
Le CM établit $f(n) = 5n + 3$ pour la somme des $n$ premiers entiers, avec
la convention : \texttt{s = s + i} vaut deux opérations, \texttt{i++} vaut
deux opérations, et il y a $n+1$ comparaisons. Un compteur incrémenté à
chaque opération transforme cette formule en fait vérifiable.
\end{lienbox}

Créez \texttt{src/main/java/Compteur.java} :

\begin{lstlisting}[caption={\texttt{Compteur.java}}]
public class Compteur {
    private long operations = 0;

    public void ajouter(int k) { operations += k; }
    public long total()        { return operations; }
    public void remise()       { operations = 0; }
}
\end{lstlisting}

Puis \texttt{src/main/java/Fragments.java} : la boucle du CM, réécrite en
déclarant chaque opération comptée.

\begin{lstlisting}[caption={\texttt{Fragments.java} --- somme instrumentée}]
public class Fragments {

    /** Somme des n premiers entiers, instrumentée. */
    public static long sommeInstrumentee(int n, Compteur c) {
        int s = 0;              c.ajouter(1);   // affectation
        int i = 0;              c.ajouter(1);   // affectation
        while (true) {
            c.ajouter(1);                       // comparaison i < n
            if (!(i < n)) break;
            s = s + i;          c.ajouter(2);   // addition + affectation
            i = i + 1;          c.ajouter(2);   // addition + affectation
        }
        return s;
    }
}
\end{lstlisting}

\begin{infobox}[Pourquoi une boucle \texttt{while} et non un \texttt{for}]
La forme \texttt{for} cache l'ordre exact des opérations : on ne voit pas
que la comparaison est évaluée une fois de plus que le corps. En dépliant
la boucle, chaque opération devient une ligne, et le compteur devient
lisible. C'est un artifice pédagogique, pas un style de production.
\end{infobox}

\begin{exercicebox}[Vérification 1 --- à compléter]
Ajoutez à \texttt{Fragments} une méthode \texttt{triangleInstrumente(int n,
Compteur c)} correspondant au fragment du CM :

\begin{center}
\texttt{for (i = 0; i < n; i++) \{ for (j = i; j < n; j++) \{ c = c + 1; \} \}}
\end{center}

N'oubliez pas de compter l'affectation \texttt{j = i} à chaque tour externe,
ni la comparaison qui fait sortir de chaque boucle.
\end{exercicebox}

% =====================================================================
\section{Étape 2 --- Un test qui force la formule}

\begin{lienbox}[Ce qu'on matérialise du CM]
Une formule de comptage est une affirmation vérifiable. Un test JUnit qui
compare le compteur à $5n+3$ pour plusieurs valeurs de $n$ transforme
l'affirmation en propriété testée --- et fait échouer le test si l'on a
oublié une opération.
\end{lienbox}

\begin{lstlisting}[caption={\texttt{src/test/java/FragmentsTest.java}}]
import org.junit.jupiter.api.Test;
import static org.junit.jupiter.api.Assertions.*;

public class FragmentsTest {

    @Test
    void sommeSuitLaFormule() {
        for (int n = 0; n <= 20; n++) {
            Compteur c = new Compteur();
            Fragments.sommeInstrumentee(n, c);
            assertEquals(5L * n + 3, c.total(), "ecart a n = " + n);
        }
    }

    @Test
    void sommeCalculeLaBonneValeur() {
        Compteur c = new Compteur();
        assertEquals(45, Fragments.sommeInstrumentee(10, c));  // 0+1+...+9
    }
}
\end{lstlisting}

\begin{lstlisting}[style=bash]
gradle test --no-daemon
\end{lstlisting}

\begin{exercicebox}[Vérification 2 --- à compléter]
\begin{enumerate}[nosep]
    \item Supprimez le \texttt{c.ajouter(1)} de la comparaison, relancez les
          tests. Quel message obtenez-vous, et pour quelle valeur de $n$ le
          test échoue-t-il en premier ? \lig{6cm}
    \item Rétablissez-le, puis écrivez le test correspondant à
          \texttt{triangleInstrumente} : la formule attendue est
          $f(n) = 2{,}5n^2 + 7{,}5n + 3$. Comment l'écrire en arithmétique
          entière, sans nombres à virgule ? \lig{6cm}
\end{enumerate}
\textit{Aide : $2{,}5n^2 + 7{,}5n + 3 = (5n^2 + 15n + 6)/2$, et le numérateur
est toujours pair.}
\end{exercicebox}

% =====================================================================
\section{Étape 3 --- Mesurer un temps d'exécution}

\begin{lienbox}[Ce qu'on matérialise du CM]
Le comptage donne $f(n)$ ; le chronomètre donne un temps. Les deux doivent
raconter la même histoire : si $f$ est quadratique, doubler $n$ doit
multiplier le temps par environ quatre. Ce rapport est l'observable qui
relie la théorie à la machine.
\end{lienbox}

\begin{lstlisting}[caption={\texttt{Benchmark.java} --- insertions en tête}]
import java.util.ArrayList;
import java.util.LinkedList;
import java.util.List;

public class Benchmark {

    public static long mesurerInsertionsEnTete(List<Integer> liste, int n) {
        long debut = System.nanoTime();
        for (int i = 0; i < n; i++) {
            liste.add(0, i);
        }
        long fin = System.nanoTime();
        return (fin - debut) / 1_000_000;   // millisecondes
    }

    public static void main(String[] args) {
        int[] tailles = {1_000, 2_000, 4_000, 8_000, 16_000, 32_000};

        // Chauffe : la JVM compile le code apres quelques milliers d'executions.
        for (int r = 0; r < 3; r++) mesurerInsertionsEnTete(new ArrayList<>(), 2_000);

        System.out.printf("%8s %14s %14s %10s%n", "n", "ArrayList(ms)", "LinkedList(ms)", "rapport");
        long precedent = 0;
        for (int n : tailles) {
            long tArray  = mesurerInsertionsEnTete(new ArrayList<>(), n);
            long tLinked = mesurerInsertionsEnTete(new LinkedList<>(), n);
            String rapport = (precedent == 0) ? "-" : String.format("%.1f", (double) tArray / precedent);
            System.out.printf("%8d %14d %14d %10s%n", n, tArray, tLinked, rapport);
            precedent = tArray;
        }
    }
}
\end{lstlisting}

\begin{lstlisting}[style=bash]
gradle run --no-daemon
\end{lstlisting}

\begin{attentionbox}[Trois pièges de la mesure]
\textbf{La chauffe :} les premières exécutions sont interprétées, les
suivantes compilées. Sans boucle de chauffe, la première ligne du tableau
est systématiquement fausse.\\
\textbf{La granularité :} en dessous de quelques millisecondes, le bruit
domine. Si une ligne affiche $0$, augmentez $n$.\\
\textbf{L'élimination de code mort :} si le résultat d'un calcul n'est
jamais utilisé, la JVM peut supprimer le calcul. Ici, la liste est bien
construite, donc le travail a lieu.
\end{attentionbox}

% =====================================================================
\section{Étape 4 --- Observer}

\begin{lienbox}[Ce qu'on matérialise du CM]
Le CM affirme que $n$ insertions en tête coûtent $n(n-1)/2$ déplacements
dans un tableau contigu et $n$ rebranchements dans une liste chaînée. Ce
programme affiche le rapport entre deux tailles consécutives : un rapport
proche de $4$ quand on double $n$ signe un coût quadratique ; un rapport
proche de $2$, un coût linéaire.
\end{lienbox}

\begin{exercicebox}[Vérification 3 --- notez vos observations]
Remplissez avec vos propres mesures.

\vspace{0.2cm}
\begin{center}
\begin{tabular}{cccc}
\toprule
$n$ & \textbf{ArrayList (ms)} & \textbf{LinkedList (ms)} & \textbf{rapport ArrayList} \\
\midrule
$1\,000$  & \rule{0pt}{0.7cm} & & --- \\
$2\,000$  & \rule{0pt}{0.7cm} & & \\
$4\,000$  & \rule{0pt}{0.7cm} & & \\
$8\,000$  & \rule{0pt}{0.7cm} & & \\
$16\,000$ & \rule{0pt}{0.7cm} & & \\
$32\,000$ & \rule{0pt}{0.7cm} & & \\
\bottomrule
\end{tabular}
\end{center}

\vspace{0.3cm}
\begin{enumerate}[nosep]
    \item Quand $n$ double, par combien le temps de l'\texttt{ArrayList}
          est-il multiplié ? \lig{3cm} \quad Cela correspond-il à la
          prédiction du CM ? \lig{4cm}
    \item Et pour la \texttt{LinkedList} ? \lig{6cm}
    \item À partir de quelle taille l'écart devient-il visible à l'œil nu ?
          \lig{4cm}
    \item Vos rapports sont-ils exactement égaux à la prédiction ? Citez une
          cause possible d'écart. \lig{7cm}
\end{enumerate}
\end{exercicebox}

\begin{exercicebox}[Vérification 4 --- une seconde mesure]
Écrivez une méthode \texttt{mesurerAccesParRang(List<Integer> liste, int n)}
qui, sur une liste déjà remplie de $n$ éléments, exécute $n$ appels à
\texttt{liste.get(i)} avec $i$ tiré au hasard, et mesure le temps total.

Prédisez le résultat \emph{avant} d'exécuter :\\
\texttt{ArrayList} : \lig{5cm} \quad \texttt{LinkedList} : \lig{5cm}

Puis exécutez. Votre prédiction était-elle correcte ? \lig{6cm}
\end{exercicebox}

% =====================================================================
\section{Enregistrer le travail}

\begin{lstlisting}[style=bash,caption={Git}]
cd /workspace/lab-s02
git add .
git commit -m "Lab S02 : instrumentation, tests de comptage, benchmarks"
\end{lstlisting}

\newpage
% #####################################################################
\part{Théorie : exercices d'application}
% #####################################################################

\begin{infobox}[Comment aborder cette partie]
Ces exercices ne sont pas un contrôle et ne sont pas notés. Ils portent sur
les fragments que vous venez d'instrumenter : le but est de retrouver par le
calcul les nombres que votre compteur a affichés. Sauf mention contraire, on
applique la convention du cours : \texttt{s = s + i} compte pour deux
opérations, \texttt{i++} pour deux, une comparaison pour une, un accès
\texttt{tab[i]} pour un.
\end{infobox}

% ---------------------------------------------------------------------
\section{Exercice 1 --- Établir $f(n)$}

Pour chaque fragment, donnez le nombre exact d'opérations puis la classe
$\bigTheta$.

\vspace{0.3cm}
\noindent\textbf{(a)}
\begin{center}
\texttt{int s = 0; for (int i = 0; i < n; i++) \{ s = s + tab[i]; \}}
\end{center}

\noindent Détail du décompte : \lig{12cm}

\noindent $f(n) = $ \lig{5cm} \qquad classe : \lig{3cm}

\vspace{0.4cm}
\noindent\textbf{(b)}
\begin{center}
\texttt{int c = 0; for (int i = 0; i < n; i++) \{ for (int j = 0; j < n; j++) \{ c = c + 1; \} \}}
\end{center}

\noindent $f(n) = $ \lig{5cm} \qquad classe : \lig{3cm}

\vspace{0.4cm}
\noindent\textbf{(c)}
\begin{center}
\texttt{int t = 1; for (int i = 1; i < n; i = i * 3) \{ t = t + i; \}}
\end{center}

\noindent Nombre de tours de boucle : \lig{4cm} \qquad classe : \lig{3cm}

\vspace{0.3cm}
\noindent Pour $n = 100$, listez les valeurs prises par \texttt{i} :
\lig{8cm}

\vspace{0.4cm}
\noindent\textbf{(d)} Comparez (b) et le fragment triangulaire du CM
(\texttt{j = i} au lieu de \texttt{j = 0}). Les deux ont-ils la même classe ?
Le même $f(n)$ ? \lig{9cm}

\vspace{0.2cm}
\noindent\lig{14cm}

% ---------------------------------------------------------------------
\section{Exercice 2 --- Prouver une borne}

\noindent\textbf{(a)} Montrez que $f(n) = 4n + 25 \in \bigO(n)$ en exhibant
un couple $(c, n_0)$ et en rédigeant l'inégalité.

\vspace{0.2cm}
\noindent\lig{14cm}

\vspace{0.2cm}
\noindent\lig{14cm}

\vspace{0.3cm}
\noindent\textbf{(b)} Le CM démontre $3n^2 + 20n + 50 \in \bigO(n^2)$ avec
$(73, 1)$, et la figure exhibe $(5, 13)$. Vérifiez que $(5, 13)$ convient
en testant $n = 12$ puis $n = 13$ :

\vspace{0.2cm}
$n = 12$ : $f(12) = $ \lig{2.5cm} \quad $5 \times 12^2 = $ \lig{2.5cm}
\quad inégalité vraie ? \lig{2cm}

\vspace{0.2cm}
$n = 13$ : $f(13) = $ \lig{2.5cm} \quad $5 \times 13^2 = $ \lig{2.5cm}
\quad inégalité vraie ? \lig{2cm}

\vspace{0.3cm}
\noindent\textbf{(c)} Peut-on conclure de (b) que $13$ est le plus petit
seuil possible \emph{quel que soit} $c$ ? \lig{7cm}

\vspace{0.2cm}
\noindent\lig{14cm}

% ---------------------------------------------------------------------
\section{Exercice 3 --- Vrai ou faux}

Cochez, sans justification.

\vspace{0.3cm}
\begin{center}
\begin{tabular}{lcc}
\toprule
\textbf{Affirmation} & \textbf{Vrai} & \textbf{Faux} \\
\midrule
$n^2 + n \in \bigO(n^2)$                      & \rule{0pt}{0.7cm} & \\
$n^2 + n \in \bigO(n^3)$                      & \rule{0pt}{0.7cm} & \\
$n^2 + n \in \bigTheta(n^3)$                  & \rule{0pt}{0.7cm} & \\
$n \log n \in \bigOmega(n)$                   & \rule{0pt}{0.7cm} & \\
$3^n \in \bigO(2^n)$                          & \rule{0pt}{0.7cm} & \\
$\bigO$ est synonyme de « pire cas »          & \rule{0pt}{0.7cm} & \\
\bottomrule
\end{tabular}
\end{center}

\vspace{0.3cm}
\noindent Pour la dernière ligne, expliquez votre réponse en une phrase :
\lig{10cm}

% ---------------------------------------------------------------------
\section{Exercice 4 --- Analyse de cas}

On considère la recherche linéaire d'une valeur dans un tableau de $n$
éléments.

\vspace{0.3cm}
\begin{center}
\begin{tabular}{lll}
\toprule
\textbf{Cas} & \textbf{Entrée qui le réalise} & \textbf{Classe} \\
\midrule
Meilleur & \rule{0pt}{0.8cm} & \\
Pire     & \rule{0pt}{0.8cm} & \\
Moyen    & \rule{0pt}{0.8cm} & \\
\bottomrule
\end{tabular}
\end{center}

\vspace{0.3cm}
\noindent\textbf{(a)} Pourquoi le pire cas sert-il de référence plutôt que
le cas moyen ? Donnez deux raisons. \lig{9cm}

\vspace{0.2cm}
\noindent\lig{14cm}

\vspace{0.3cm}
\noindent\textbf{(b)} Décrivez l'entrée réalisant le pire cas pour la
recherche d'un doublon par double boucle avec sortie anticipée :
\lig{8cm}

% ---------------------------------------------------------------------
\section{Exercice 5 --- Récurrences}

Résolvez par le théorème principal. Pour chacune, donnez $a$, $b$,
$\alpha = \log_b a$, le cas applicable et la conclusion.

\vspace{0.3cm}
\begin{center}
\begin{tabular}{lcccc}
\toprule
\textbf{Récurrence} & $a$ & $b$ & $\alpha$ & \textbf{Cas et conclusion} \\
\midrule
$T(n) = T(n/2) + \bigTheta(1)$      & \rule{0pt}{0.8cm} & & & \\
$T(n) = 2T(n/2) + \bigTheta(1)$     & \rule{0pt}{0.8cm} & & & \\
$T(n) = 2T(n/2) + \bigTheta(n)$     & \rule{0pt}{0.8cm} & & & \\
$T(n) = 3T(n/2) + \bigTheta(n^2)$   & \rule{0pt}{0.8cm} & & & \\
\bottomrule
\end{tabular}
\end{center}

\vspace{0.3cm}
\noindent\textbf{(a)} Les deux premières lignes ne diffèrent que par $a$.
Expliquez, sans calcul, pourquoi l'écart de conclusion est si grand :
\lig{9cm}

\vspace{0.2cm}
\noindent\lig{14cm}

\vspace{0.3cm}
\noindent\textbf{(b)} À laquelle de ces récurrences correspond la recherche
dichotomique ? \lig{5cm}

% ---------------------------------------------------------------------
\section{Exercice 6 --- Question subsidiaire}

Vos mesures de la Partie I ne donnent probablement pas un rapport exactement
égal à $4$ quand $n$ double.

\vspace{0.2cm}
\noindent\textbf{(a)} Citez deux causes possibles, l'une liée à la machine,
l'autre à la théorie elle-même : \lig{9cm}

\vspace{0.2cm}
\noindent\lig{14cm}

\vspace{0.3cm}
\noindent\textbf{(b)} Un écart mesuré invalide-t-il la prédiction
asymptotique ? Argumentez en deux phrases. \lig{9cm}

\vspace{0.2cm}
\noindent\lig{14cm}

\newpage
% =====================================================================
\section*{Pour aller plus loin (optionnel)}
\addcontentsline{toc}{section}{Pour aller plus loin (optionnel)}

\begin{itemize}[nosep]
    \item Tracer les courbes mesurées avec matplotlib (le conteneur dispose
          de Python et matplotlib) et superposer $n^2/2$ ajusté par un
          facteur constant. La forme doit coïncider, pas les valeurs.
    \item Instrumenter votre \texttt{ArrayDeque} du Lab~S01 : compter les
          déplacements réellement effectués par \texttt{agrandir()} sur une
          séquence de $n$ insertions, et comparer à $n$.
    \item Reprendre la file naïve sur \texttt{ArrayList} avec
          \texttt{remove(0)} et mesurer $n$ \texttt{dequeue} successifs.
\end{itemize}

% =====================================================================
\section*{Aide-mémoire}
\addcontentsline{toc}{section}{Aide-mémoire}

\begin{center}
\begin{tabular}{ll}
\toprule
\textbf{Question} & \textbf{Réponse} \\
\midrule
\texttt{s = s + i} & 2 opérations \\
Comparaisons d'une boucle de $n$ tours & $n + 1$ \\
Boucles successives & on additionne, le maximum l'emporte \\
Boucles imbriquées & on multiplie \\
Pas multiplicatif (\texttt{i *= k}) & $\bigTheta(\log n)$ tours \\
Prouver $f \in \bigO(g)$ & exhiber $c > 0$ et $n_0$ \\
Rapport de temps quand $n$ double & $\approx 2$ si linéaire, $\approx 4$ si quadratique \\
$T(n) = aT(n/b) + f(n)$ & comparer $f(n)$ à $n^{\log_b a}$ \\
\midrule
Compiler / tester / exécuter & \texttt{gradle build | test | run -\/-no-daemon} \\
\bottomrule
\end{tabular}
\end{center}

% =====================================================================
\begin{attentionbox}[Rapport individuel --- obligatoire]
Chaque étudiant rédige un rapport de \textbf{1 à 2 pages} :
\begin{enumerate}[nosep]
    \item \textbf{Ce que j'ai fait :} fragments instrumentés, tests écrits,
          mesures obtenues.
    \item \textbf{Les observations clés :} le tableau de mesures
          confirme-t-il la prédiction du CM ? Où l'écart apparaît-il ?
    \item \textbf{Les difficultés :} formules de comptage fausses, mesures
          instables, blocages sur les exercices théoriques.
    \item \textbf{Ce que j'ai appris :} le concept le plus important retenu
          de cette séance.
\end{enumerate}
\textbf{Format :} PDF. \textbf{Nom :} \texttt{rapport-lab-S02-PRENOM-NOM.pdf}
\end{attentionbox}

\end{document}
